\documentclass[11pt,a4paper]{article}
\usepackage{xcolor}

\definecolor{comentcolor}{HTML}{0070C0}

\newcommand{\coment}[1]{%
    \textcolor{comentcolor}{\textbf{#1}}%
}

% =========================================================
% PAQUETES
% =========================================================
\usepackage[utf8]{inputenc}
\usepackage[T1]{fontenc}
\usepackage[spanish,es-nodecimaldot]{babel}
\usepackage{amsmath}
\usepackage{amsfonts}
\usepackage{amssymb}
\usepackage{graphicx}
\usepackage{geometry}
\usepackage{enumitem}
\usepackage{titlesec}

% =========================================================
% CONFIGURACIÓN DE MÁRGENES
% =========================================================
\geometry{
    top=2.5cm,
    bottom=2.5cm,
    left=2.5cm,
    right=2.5cm
}

% =========================================================
% CONFIGURACIÓN DE TÍTULOS
% =========================================================
\renewcommand{\thesection}{\Roman{section}}

\titleformat{\section}
    {\large\bfseries}
    {\thesection.}
    {0.5em}
    {\MakeUppercase}

% =========================================================
% CONFIGURACIÓN DE LISTAS
% =========================================================
\setlist[enumerate,1]{
    label=\textbf{\arabic*.},
    leftmargin=*,
    itemsep=1em,
    topsep=0.5em
}

\setlist[enumerate,2]{
    label=\alph*),
    leftmargin=1.5cm,
    itemsep=0.3em
}

\setlist[itemize]{
    label=--,
    leftmargin=1.5cm,
    itemsep=0.2em
}

% =========================================================
% INFORMACIÓN DEL DOCUMENTO
% =========================================================
\title{\textbf{Ejercicios}}
\author{Modelo sistema: Matriz Inversa / Regla de Cramer}
\date{}

% =========================================================
% INICIO DEL DOCUMENTO
% =========================================================
\begin{document}

\maketitle

% =========================================================
% SECCIÓN I
% =========================================================


\begin{enumerate}

\item Una plataforma digital procesa solicitudes mediante tres operaciones: \emph{análisis con inteligencia artificial}, \emph{verificación de identidad} y \emph{almacenamiento cifrado}. Cada operación consume una cantidad fija de créditos computacionales por solicitud. Se quiere determinar el consumo individual de estas operaciones para optimizar el servicio.

Una solicitud que utiliza las tres operaciones requiere {95 créditos}. Cuando el análisis con inteligencia artificial no se ejecuta, la misma solicitud requiere {55 créditos}. Por último, en esa solicitud reducida se reemplaza el almacenamiento cifrado por el análisis con inteligencia artificial, manteniendo la verificación de identidad; el consumo total {aumenta 15 créditos} respecto de los 55 créditos anteriores.




\begin{enumerate}
\item Defina las tres variables, plantee el sistema de ecuaciones lineales que modela la situación y escriba su representación matricial $Ax=b$, identificando $A$, $x$ y $b$.
\item Resuelva el sistema mediante el método de la matriz inversa y determine los valores de las tres incógnitas en el contexto del problema.
\end{enumerate}

\item  Una empresa de envases reutilizables cobra a sus clientes por tres conceptos: \emph{reacondicionamiento de cada envase}, \emph{servicio de entrega} y \emph{depósito reembolsable}. Este último se devuelve íntegramente cuando el cliente retorna el envase. La empresa quiere identificar cuánto corresponde a cada concepto, suponiendo precios constantes.

Una clienta paga inicialmente {\$95} por reacondicionar un envase, recibir una entrega y dejar el depósito. Cuando devuelve el envase, su gasto definitivo es de {\$65}. Por otra parte, un cliente lleva {dos envases propios} para reacondicionarlos y solicita {una sola entrega}. Como no necesita dejar depósito, paga un total de {\$105}.




\begin{enumerate}
\item Defina las tres variables, plantee el sistema de ecuaciones lineales que modela la situación y escriba su representación matricial $Ax=b$, identificando $A$, $x$ y $b$.
\item Resuelva el sistema mediante el método de la matriz inversa y determine los valores de las tres incógnitas en el contexto del problema.
\end{enumerate}

\item Un centro de datos ejecuta simultáneamente tres tipos de
procesos informáticos:

\begin{itemize}
\item Procesos de análisis de datos (A),
\item Procesos de respaldo y almacenamiento de información (B),
\item Procesos de inteligencia artificial para reconocimiento de patrones (C).
\end{itemize}

Cada proceso consume recursos computacionales del servidor.
Un proceso de análisis de datos utiliza 2 unidades de CPU,
1 unidad de memoria RAM y 1 unidad de almacenamiento.
Un proceso de respaldo utiliza 3 unidades de CPU y 2 unidades
de memoria RAM. Por su parte, un proceso de inteligencia
artificial utiliza 2 unidades de CPU, 1 unidad de memoria RAM
y 2 unidades de almacenamiento. Durante una hora de operación se registró un consumo total de 21 unidades de CPU, 12 unidades de memoria RAM y
10 unidades de almacenamiento.

\begin{enumerate}
\item Defina las tres variables, plantee el sistema de ecuaciones
lineales que modela la situación y escriba su representación
matricial $Ax=b$, identificando $A$, $x$ y $b$.

\item Resuelva el sistema mediante la regla de Cramer y determine los valores de las tres incógnitas
en el contexto del problema.
\end{enumerate}

\item Una empresa constructora fabrica tres tipos de elementos
prefabricados de hormigón: vigas, columnas y losas.
Para producirlos se emplean cemento, arena y grava.

Cada viga requiere 2 unidades de cemento, 1 unidad de arena
y 1 unidad de grava. Cada columna requiere 3 unidades de
cemento y 2 unidades de arena. Cada losa requiere 2 unidades
de cemento, 1 unidad de arena y 2 unidades de grava.Durante una jornada de producción se consumieron exactamente 28 unidades de cemento, 16 unidades de arena y
13 unidades de grava.

\begin{enumerate}
\item Defina las tres variables, plantee el sistema de ecuaciones
lineales que modela la situación y escriba su representación
matricial $Ax=b$, identificando $A$, $x$ y $b$.

\item Resuelva el sistema mediante la regla de Cramer y determine los valores de las tres incógnitas
en el contexto del problema.
\end{enumerate}

\item Una fábrica produce tres modelos de sillas.
Para confeccionar una unidad del Modelo 1, se requieren
2 horas de carpintería, 1 kg de barniz y 1 hora de tapizado.
Para una silla del Modelo 2, se necesita 1 hora de carpintería,
2 kg de barniz y 1 hora de tapizado.
Por último, el Modelo 3 consume 1 hora de carpintería,
1 kg de barniz y 1 hora de tapizado. Al finalizar esta semana, los registros de producción indicaron que se emplearon exactamente 6 horas de carpintería,
6 kg de barniz y 4 horas de tapizado, sin que sobrara
ningún recurso.




\begin{enumerate}
\item Defina las tres variables, plantee el sistema de ecuaciones lineales que modela la situación y escriba su representación matricial $Ax=b$, identificando $A$, $x$ y $b$.
\item Resuelva el sistema mediante la regla de Cramer  y determine los valores de las tres incógnitas en el contexto del problema.
\end{enumerate}


\item Una planta de producción química utiliza tres reactivos,
$R_1$, $R_2$ y $R_3$, para elaborar una solución industrial
empleada en el tratamiento de aguas. Cada reactivo aporta
diferentes cantidades de tres sustancias activas denominadas
A, B y C, las cuales son necesarias para garantizar las
propiedades químicas del producto final.

Se sabe que cada litro del reactivo $R_1$ aporta
3 unidades de la sustancia A, 3 unidades de la sustancia B
y 2 unidades de la sustancia C.
Cada litro del reactivo $R_2$ aporta 2 unidades de la sustancia A,
3 unidades de la sustancia B y 2 unidades de la sustancia C.
Finalmente, cada litro del reactivo $R_3$ aporta
2 unidades de la sustancia A y 1 unidad de la sustancia C.

Después de realizar la mezcla, el laboratorio de control
de calidad determinó que la solución obtenida contiene
exactamente 9 unidades de la sustancia A,
6 unidades de la sustancia B y 6 unidades de la sustancia C.

\begin{enumerate}
\item Defina las tres variables, plantee el sistema de ecuaciones
lineales que modela la situación y escriba su representación
matricial $Ax=b$, identificando $A$, $x$ y $b$.

\item Resuelva el sistema mediante la regla de Cramer y determine los valores de las tres incógnitas
en el contexto del problema.
\end{enumerate}

\item En un laboratorio de instrumentación electrónica se
emplean tres sensores, $S_1$, $S_2$ y $S_3$, para monitorear
una variable física mediante un sistema de adquisición de
datos compuesto por tres canales de medición. La señal registrada en cada canal corresponde a la suma de las contribuciones generadas por los tres sensores.
Por cada unidad de señal emitida:

El sensor $S_1$ aporta 2 unidades al Canal 1
y 2 unidades al Canal 2. El sensor $S_2$ aporta 1 unidad al Canal 1,
2 unidades al Canal 2 y 1 unidad al Canal 3. El sensor $S_3$ aporta 1 unidad al Canal 2. Durante una prueba de calibración, el sistema registró
un total de 10 unidades de señal en el Canal 1,
16 unidades en el Canal 2 y 4 unidades en el Canal 3.

\begin{enumerate}
\item Defina las tres variables, plantee el sistema de ecuaciones
lineales que modela la situación y escriba su representación
matricial $Ax=b$, identificando $A$, $x$ y $b$.

\item Resuelva el sistema mediante el método de la matriz inversa y determine los valores de las tres
incógnitas en el contexto del problema.
\end{enumerate}

\item Se desea montar una placa triangular de acero de
espesor uniforme. La placa se va a suspender horizontalmente
desde el techo del taller mediante 3 cables verticales
ubicados en sus tres vértices A, B y C.

Se necesita calcular la tensión en cada cable para
seleccionar el diámetro correcto.

La suma de las tres tensiones debe soportar el peso total
de la placa que es de 6 kN.

Por equilibrio de momentos respecto al borde $x$,
la tensión $T_2$ aporta el doble que $T_3$ y $T_1$
no aporta momento; el momento total requerido respecto
al borde $x$ es de 7 kNm.

Por equilibrio de momentos respecto a $y$,
$T_1$ aporta 1 vez, $T_2$ aporta 2 veces y
$T_3$ aporta 2 veces; el momento total requerido
respecto a $y$ es de 10 kNm.




\begin{enumerate}
\item Defina las tres variables, plantee el sistema de ecuaciones lineales que modela la situación y escriba su representación matricial $Ax=b$, identificando $A$, $x$ y $b$.
\item Resuelva el sistema mediante la regla de Cramer y determine los valores de las tres incógnitas en el contexto del problema.
\end{enumerate}

\item Se desea preparar un fertilizante que contenga
tres tipos de abono.

\begin{itemize}
\item El abono A: no contiene Nitrógeno,
contiene 2 kg de Fósforo y 1 kg de Potasio por saco.

\item El abono B: contiene 1 kg de Nitrógeno,
1 kg de Fósforo y 1 kg de Potasio por saco.

\item El abono C: contiene 2 kg de Nitrógeno,
3 kg de Fósforo y 1 kg de Potasio por saco.
\end{itemize}

La mezcla del fertilizante requerido debe contener
exactamente 7 kg de Nitrógeno, 14 kg de Fósforo
y 6 kg de Potasio.




\begin{enumerate}
\item Defina las tres variables, plantee el sistema de ecuaciones lineales que modela la situación y escriba su representación matricial $Ax=b$, identificando $A$, $x$ y $b$.
\item Resuelva el sistema mediante el método de la matriz inversa y determine los valores de las tres incógnitas en el contexto del problema.
\end{enumerate}

\end{enumerate}

% =========================================================
% FIN DEL DOCUMENTO
% =========================================================
\end{document}